\documentclass[11pt]{article}

\usepackage[T1]{fontenc}
\usepackage[utf8]{inputenc}
\usepackage[margin=1in]{geometry}
\usepackage{amsmath,amssymb,amsthm,mathtools}
\usepackage[colorlinks=true,linkcolor=blue,citecolor=blue,urlcolor=blue]{hyperref}

\numberwithin{equation}{section}
\newtheorem{theorem}{Theorem}[section]
\newtheorem{proposition}[theorem]{Proposition}
\newtheorem{lemma}[theorem]{Lemma}
\newtheorem{corollary}[theorem]{Corollary}
\newtheorem{conjecture}[theorem]{Conjecture}
\newtheorem{problem}[theorem]{Problem}
\theoremstyle{definition}
\newtheorem{definition}[theorem]{Definition}
\newtheorem{example}[theorem]{Example}
\theoremstyle{remark}
\newtheorem{remark}[theorem]{Remark}

\DeclareMathOperator{\Var}{Var}
\DeclareMathOperator{\Prob}{P}
\DeclareMathOperator{\supp}{supp}
\newcommand{\E}{\mathbb E}
\newcommand{\N}{\mathbb N}
\newcommand{\Z}{\mathbb Z}
\newcommand{\1}{\mathbf 1}

\title{Zeta-Law Entropy, Modular Resolution, and Applications to Zeta Inequalities}
\author{/project-pudim}
\date{}

\begin{document}
\maketitle

\begin{abstract}
We study the Riemann zeta function as the partition function of the probability law \(\rho_\beta(n)=n^{-\beta}/\zeta(\beta)\) on the positive integers.  This normalization turns zeta ratios into moments, logarithmic derivatives into energy cumulants, and divisor identities into probabilistic decompositions.  The first main result identifies the microscopic successor entropy of the zeta law with the limit and supremum of finite modular successor entropies; for prime moduli these modular shadows are explicit nonlinear functionals of Dirichlet \(L\)-values.  The second structural result is an Euler-score decomposition expressing \(-(\log\zeta)'(\beta)\) as a von Mangoldt-weighted divisibility average.  We then use the same calculus to prove three zeta-inequality applications: positivity of a Nantomah expression involving \(\zeta(n+1),\zeta(n+2),\zeta(n+3)\); the Alzer--Kwong convexity and concavity pattern for \(1/\zeta\) on negative real intervals; and a generalized Holder inequality for \(\Gamma(s)\zeta(s)\) that extends Sroysang's formulation.  The framework gives a compact route from zeta probability laws to entropy, curvature, and Mellin-integral inequalities.
\end{abstract}

\noindent\textbf{Keywords.} Riemann zeta function; zeta probability law; Gibbs measure; entropy; modular residues; Dirichlet \(L\)-functions; von Mangoldt function; reciprocal zeta; Gamma-zeta inequalities; Holder inequality.

\section{Motivation and main results}

Many classical identities for \(\zeta(s)\) become more transparent after a probabilistic normalization.  For \(\beta>1\), the weights \(n^{-\beta}\) define a Gibbs law whose partition function is precisely \(\zeta(\beta)\).  The energy is \(\log n\), the free energy is \(\log\zeta(\beta)\), and ratios such as \(\zeta(\beta+r)/\zeta(\beta)\) become moments.  This viewpoint is elementary, but it packages several unrelated-looking zeta inequalities into a common calculus.

The paper develops four connected layers.  The moment layer records basic zeta-law identities.  The modular entropy layer compares the successor map \(n\mapsto n+1\) with its finite residue shadows modulo \(q\).  The Euler-score layer rewrites zeta energy through von Mangoldt divisibility events.  The Mellin-Planck layer treats \(\Gamma(s)\zeta(s)\) as a continuous partition function and converts generalized Holder inequalities into integral norm inequalities.

The application section gives three solved external problems.  The Nantomah positivity question is proved by splitting a zeta expression into a positive linear tail and a controlled quadratic correction.  The Alzer--Kwong reciprocal-zeta conjecture is proved by transporting a positive-axis curvature estimate through the functional equation.  The Sroysang generalized Holder problem is proved directly from generalized Holder applied to the Planck kernel.

\section{Notation and zeta-law calculus}

\begin{definition}[Riemann zeta probability law]
\label{def:zeta-law}
For \(\beta>1\), define a probability law \(\rho_\beta\) on \(\N\) by
\begin{equation}
\label{eq:zeta-law}
\rho_\beta(n)=\frac{n^{-\beta}}{\zeta(\beta)},\qquad n\in\N.
\end{equation}
Equivalently, with \(E(n)=\log n\) and \(Z(\beta)=\zeta(\beta)\),
\begin{equation}
\label{eq:gibbs-form}
\rho_\beta(n)=\frac{e^{-\beta E(n)}}{Z(\beta)}.
\end{equation}
\end{definition}

\begin{definition}[Zeta free energy]
\label{def:zeta-free-energy}
The zeta free energy is
\begin{equation}
\label{eq:zeta-free-energy}
A(\beta)=\log\zeta(\beta),\qquad \beta>1.
\end{equation}
\end{definition}

\begin{proposition}[Free-energy derivatives]
\label{prop:free-energy-derivatives}
For every \(\beta>1\),
\begin{align}
\label{eq:free-energy-derivatives}
A'(\beta)&=-\E_\beta[\log N],&
A''(\beta)&=\Var_\beta(\log N).
\end{align}
\end{proposition}

\begin{proof}
Absolute convergence of the differentiated Dirichlet series on compact subintervals of \((1,\infty)\) permits termwise differentiation.  Thus
\begin{equation}
\label{eq:first-free-derivative-proof}
A'(\beta)=\frac{\zeta'(\beta)}{\zeta(\beta)}
=-\frac{\sum_{n\ge1}(\log n)n^{-\beta}}{\zeta(\beta)}
=-\E_\beta[\log N].
\end{equation}
Differentiating once more gives
\begin{equation}
\label{eq:second-free-derivative-proof}
A''(\beta)=\E_\beta[(\log N)^2]-\E_\beta[\log N]^2
=\Var_\beta(\log N).
\end{equation}
\end{proof}

\begin{proposition}[Zeta-law calculus]
\label{prop:zeta-law-calculus}
Let \(\beta>1\), and let \(N\) have law \(\rho_\beta\).  For every \(r\ge0\),
\begin{equation}
\label{eq:zeta-law-moment}
\E_\beta[N^{-r}]=\frac{\zeta(\beta+r)}{\zeta(\beta)}.
\end{equation}
Moreover, for \(\alpha,\beta>1\),
\begin{equation}
\label{eq:zeta-law-relative-entropy}
D(\rho_\alpha\Vert\rho_\beta)
=(\beta-\alpha)\E_\alpha[\log N]
+\log\zeta(\beta)-\log\zeta(\alpha).
\end{equation}
\end{proposition}

\begin{proof}
The moment identity follows by direct summation:
\begin{equation}
\label{eq:zeta-law-moment-proof}
\E_\beta[N^{-r}]
=\sum_{n=1}^{\infty}n^{-r}\frac{n^{-\beta}}{\zeta(\beta)}
=\frac{\zeta(\beta+r)}{\zeta(\beta)}.
\end{equation}
For relative entropy, substitute \eqref{eq:zeta-law}:
\begin{equation}
\label{eq:zeta-law-relative-entropy-proof}
D(\rho_\alpha\Vert\rho_\beta)
=\sum_{n=1}^{\infty}\rho_\alpha(n)
\left((\beta-\alpha)\log n+\log\zeta(\beta)-\log\zeta(\alpha)\right),
\end{equation}
which is \eqref{eq:zeta-law-relative-entropy}.
\end{proof}

\section{Modular successor entropy}

\begin{definition}[Modular residue distribution and successor entropy]
\label{def:modular-residue-successor}
For \(q\ge2\), define the residue distribution of \(\rho_\beta\) modulo \(q\) by
\begin{equation}
\label{eq:modular-residue-distribution}
\mu_{q,\beta}(a)=
\sum_{\substack{n\ge1\\ n\equiv a\pmod q}}\rho_\beta(n),
\qquad a\in\Z/q\Z.
\end{equation}
Writing \(a+1\) cyclically in \(\Z/q\Z\), define
\begin{equation}
\label{eq:modular-successor-entropy}
B_q(\beta)=
\sum_{a\in\Z/q\Z}
\mu_{q,\beta}(a)\log\frac{\mu_{q,\beta}(a)}{\mu_{q,\beta}(a+1)}.
\end{equation}
\end{definition}

\begin{theorem}[Zeta-law successor entropy and modular resolution]
\label{thm:successor-entropy}
For every \(\beta>1\), define
\begin{equation}
\label{eq:successor-entropy}
\Sigma(\beta)=\sum_{n=1}^{\infty}\rho_\beta(n)
\log\frac{\rho_\beta(n)}{\rho_\beta(n+1)}.
\end{equation}
Then
\begin{equation}
\label{eq:successor-entropy-series}
\Sigma(\beta)=
\frac{\beta}{\zeta(\beta)}
\sum_{k=1}^{\infty}\frac{(-1)^{k+1}}{k}\zeta(\beta+k),
\end{equation}
and
\begin{equation}
\label{eq:successor-entropy-modular-resolution}
\Sigma(\beta)=\sup_{q\ge2}B_q(\beta)=\lim_{q\to\infty}B_q(\beta).
\end{equation}
\end{theorem}

\begin{proof}
Since
\begin{equation}
\label{eq:successor-ratio}
\frac{\rho_\beta(n)}{\rho_\beta(n+1)}
=\left(1+\frac1n\right)^\beta,
\end{equation}
we have
\begin{equation}
\label{eq:successor-entropy-log-sum}
\Sigma(\beta)=\frac{\beta}{\zeta(\beta)}
\sum_{n=1}^{\infty}n^{-\beta}\log\left(1+\frac1n\right).
\end{equation}
Using
\begin{equation}
\label{eq:log-series}
\log(1+x)=\sum_{k=1}^{\infty}\frac{(-1)^{k+1}}{k}x^k
\qquad(0<x\le1)
\end{equation}
with Abel limiting at \(x=1\), and summing against \(n^{-\beta}\), gives \eqref{eq:successor-entropy-series}.

For the modular identity, fix \(q\ge2\).  The log-sum inequality gives, for each residue class \(a\),
\begin{equation}
\label{eq:modular-log-sum}
\sum_{\substack{n\ge1\\ n\equiv a\pmod q}}
\rho_\beta(n)\log\frac{\rho_\beta(n)}{\rho_\beta(n+1)}
\ge
\mu_{q,\beta}(a)
\log\frac{\mu_{q,\beta}(a)}
{\sum_{n\equiv a\pmod q}\rho_\beta(n+1)}.
\end{equation}
For \(a\ne0\), the denominator is \(\mu_{q,\beta}(a+1)\).  For \(a=0\), it is \(\mu_{q,\beta}(1)-\rho_\beta(1)\le\mu_{q,\beta}(1)\).  Summing \eqref{eq:modular-log-sum} over \(a\) gives
\begin{equation}
\label{eq:successor-dominates-modular}
\Sigma(\beta)\ge B_q(\beta)
\qquad(q\ge2).
\end{equation}

Conversely, fix \(M\in\N\).  If \(q>M+1\), then for \(1\le a\le M+1\),
\begin{equation}
\label{eq:residue-isolation}
\mu_{q,\beta}(a)=\rho_\beta(a)+O_{\beta,M}(q^{-\beta}),
\end{equation}
and \(\mu_{q,\beta}(0)=q^{-\beta}\).  All cyclic terms in \(B_q(\beta)\) are nonnegative except possibly the boundary term, which satisfies
\begin{equation}
\label{eq:boundary-term}
\mu_{q,\beta}(0)\log\frac{\mu_{q,\beta}(0)}{\mu_{q,\beta}(1)}
=O_\beta(q^{-\beta}\log q).
\end{equation}
Therefore
\begin{equation}
\label{eq:modular-liminf}
\liminf_{q\to\infty}B_q(\beta)
\ge
\sum_{n=1}^{M}\rho_\beta(n)\log\frac{\rho_\beta(n)}{\rho_\beta(n+1)}.
\end{equation}
Letting \(M\to\infty\) gives \(\liminf_qB_q(\beta)\ge\Sigma(\beta)\), which together with \eqref{eq:successor-dominates-modular} proves \eqref{eq:successor-entropy-modular-resolution}.
\end{proof}

\begin{corollary}[Prime-modulus Dirichlet L-resolution]
\label{cor:prime-dirichlet-resolution}
Let \(p\) be prime, let \(\chi\) range over Dirichlet characters modulo \(p\), and define
\begin{equation}
\label{eq:prime-dirichlet-shadow}
S_a(p,\beta)=\sum_{\chi\bmod p}\overline{\chi(a)}L(\beta,\chi)
\qquad(1\le a\le p-1).
\end{equation}
Then
\begin{equation}
\label{eq:prime-residue-formula}
\mu_{p,\beta}(a)=\frac{S_a(p,\beta)}{(p-1)\zeta(\beta)}
\quad(1\le a\le p-1),
\qquad
\mu_{p,\beta}(0)=p^{-\beta}.
\end{equation}
Consequently the prime-modulus functionals obtained by substituting \eqref{eq:prime-residue-formula} into \(B_p(\beta)\) satisfy
\begin{equation}
\label{eq:prime-functional-limit}
\Sigma(\beta)=
\lim_{\substack{p\to\infty\\ p\ {\rm prime}}}
F_p\bigl(\zeta(\beta),\{L(\beta,\chi)\}_{\chi\bmod p}\bigr)
=
\sup_{p\ {\rm prime}}
F_p\bigl(\zeta(\beta),\{L(\beta,\chi)\}_{\chi\bmod p}\bigr).
\end{equation}
\end{corollary}

\begin{proof}
For \(a\not\equiv0\pmod p\), character orthogonality gives
\begin{equation}
\label{eq:character-orthogonality}
\frac1{p-1}\sum_{\chi\bmod p}\overline{\chi(a)}\chi(n)
=
\begin{cases}
1,& n\equiv a\pmod p,\\
0,& n\not\equiv a\pmod p.
\end{cases}
\end{equation}
Therefore
\begin{equation}
\label{eq:dirichlet-residue-sum}
\sum_{\substack{n\ge1\\ n\equiv a\pmod p}}n^{-\beta}
=\frac1{p-1}\sum_{\chi\bmod p}\overline{\chi(a)}L(\beta,\chi).
\end{equation}
Dividing by \(\zeta(\beta)\) gives the nonzero residue formula, while
\begin{equation}
\label{eq:zero-residue-prime}
\mu_{p,\beta}(0)=\frac1{\zeta(\beta)}\sum_{m=1}^{\infty}(mp)^{-\beta}
=p^{-\beta}.
\end{equation}
The functional identity follows by substituting these residue formulas into \(B_p(\beta)\) and applying Theorem \ref{thm:successor-entropy} along the prime sequence.
\end{proof}

\section{Euler-score identities}

\begin{theorem}[Euler-score decomposition]
\label{thm:euler-score-decomposition}
For every \(\beta>1\),
\begin{equation}
\label{eq:euler-score-decomposition}
-\bigl(\log\zeta\bigr)'(\beta)
=\E_\beta[\log N]
=\sum_{d=1}^{\infty}\Lambda(d)\mu_{d,\beta}(0)
=\sum_{d=1}^{\infty}\frac{\Lambda(d)}{d^\beta}.
\end{equation}
Equivalently,
\begin{equation}
\label{eq:log-zeta-von-mangoldt}
-\frac{\zeta'(\beta)}{\zeta(\beta)}
=\sum_{d=1}^{\infty}\frac{\Lambda(d)}{d^\beta}.
\end{equation}
\end{theorem}

\begin{proof}
The von Mangoldt identity is
\begin{equation}
\label{eq:mangoldt-identity}
\log n=\sum_{d\mid n}\Lambda(d).
\end{equation}
By nonnegative summation,
\begin{equation}
\label{eq:euler-score-proof}
\E_\beta[\log N]
=\sum_{n\ge1}\rho_\beta(n)\sum_{d\mid n}\Lambda(d)
=\sum_{d\ge1}\Lambda(d)\Prob_\beta(d\mid N).
\end{equation}
The divisibility probability is
\begin{equation}
\label{eq:divisibility-probability}
\Prob_\beta(d\mid N)
=\frac1{\zeta(\beta)}\sum_{m\ge1}(dm)^{-\beta}
=d^{-\beta}.
\end{equation}
Substitution into \eqref{eq:euler-score-proof} proves \eqref{eq:euler-score-decomposition}; \eqref{eq:log-zeta-von-mangoldt} follows from \eqref{eq:first-free-derivative-proof}.
\end{proof}

\begin{proposition}[Finite Euler-score identity]
\label{prop:finite-euler-score}
For every \(A\subseteq\{1,\ldots,N\}\),
\begin{equation}
\label{eq:finite-euler-score}
\sum_{n\in A}\log n
=\sum_{d\le N}\Lambda(d)\sum_{m\le N/d}\1_A(md).
\end{equation}
\end{proposition}

\begin{proof}
Reverse the order of summation:
\begin{equation}
\label{eq:finite-euler-score-proof}
\sum_{d\le N}\Lambda(d)\sum_{m\le N/d}\1_A(md)
=\sum_{n\in A}\sum_{d\mid n}\Lambda(d)
=\sum_{n\in A}\log n.
\end{equation}
\end{proof}

\section{Curvature and Mellin-Planck tools}

\begin{lemma}[Uniform positive-axis curvature bound]
\label{lem:uniform-positive-axis-curvature}
For every \(s\ge3\),
\begin{equation}
\label{eq:uniform-curvature-bound}
(\log\zeta)''(s)=\Var_s(\log N)<\frac13.
\end{equation}
\end{lemma}

\begin{proof}
By Proposition \ref{prop:free-energy-derivatives},
\begin{equation}
\label{eq:curvature-first-bound}
(\log\zeta)''(s)
=\frac{\zeta''(s)}{\zeta(s)}
-\left(\frac{\zeta'(s)}{\zeta(s)}\right)^2
\le \frac{\zeta''(s)}{\zeta(s)}
\le \zeta''(s).
\end{equation}
For \(s\ge3\),
\begin{equation}
\label{eq:zeta-second-derivative-bound}
\zeta''(s)=\sum_{m=2}^{\infty}\frac{(\log m)^2}{m^s}
\le\sum_{m=2}^{\infty}\frac{(\log m)^2}{m^3}.
\end{equation}
The function \(x\mapsto(\log x)^2x^{-3}\) is decreasing for \(x\ge2\), so
\begin{equation}
\label{eq:curvature-integral-comparison}
\sum_{m=2}^{\infty}\frac{(\log m)^2}{m^3}
\le
\frac{(\log2)^2}{8}
+\int_2^\infty\frac{(\log x)^2}{x^3}\,dx.
\end{equation}
Since
\begin{equation}
\label{eq:curvature-integral-value}
\int_2^\infty\frac{(\log x)^2}{x^3}\,dx
=\frac{(\log2)^2+\log2+\frac12}{8},
\end{equation}
we obtain
\begin{equation}
\label{eq:curvature-final-bound}
(\log\zeta)''(s)
\le
\frac{2(\log2)^2+\log2+\frac12}{8}
<\frac13.
\end{equation}
\end{proof}

\begin{definition}[Mellin-Planck partition function]
\label{def:mellin-planck-partition}
For \(s>1\), define
\begin{equation}
\label{eq:mellin-planck-partition}
M(s)=\Gamma(s)\zeta(s)
=\int_0^\infty\frac{t^{s-1}}{e^t-1}\,dt.
\end{equation}
The corresponding probability law is
\begin{equation}
\label{eq:mellin-planck-law}
\nu_s(dt)=\frac{t^{s-1}}{(e^t-1)M(s)}\,dt.
\end{equation}
\end{definition}

\begin{proposition}[Mellin-Planck log-convexity]
\label{prop:mellin-planck-log-convexity}
For \(s>1\),
\begin{equation}
\label{eq:mellin-planck-log-convexity}
\frac{d^2}{ds^2}\log M(s)=\Var_{\nu_s}(\log t)\ge0.
\end{equation}
\end{proposition}

\begin{proof}
Differentiate the logarithm of the integral in \eqref{eq:mellin-planck-partition} under the integral sign.  The first derivative is the \(\nu_s\)-expectation of \(\log t\), and the second derivative is its variance:
\begin{equation}
\label{eq:mellin-planck-variance-proof}
\frac{d^2}{ds^2}\log M(s)
=\E_{\nu_s}[(\log t)^2]-\E_{\nu_s}[\log t]^2.
\end{equation}
\end{proof}

\section{Applications to external zeta inequalities}

\begin{theorem}[Affirmative solution of Nantomah zeta positivity problem]
\label{thm:nantomah-positivity}
For every \(n\in\N\),
\begin{equation}
\label{eq:nantomah-positivity}
(n+2)\zeta(n+1)\zeta(n+3)
-(n+1)\zeta(n+2)^2
-\zeta(n+1)\zeta(n+2)>0.
\end{equation}
\end{theorem}

\begin{proof}
Put \(s=n+1\ge2\) and \(Z_s=\zeta(s)\).  The expression in \eqref{eq:nantomah-positivity} becomes
\begin{equation}
\label{eq:nantomah-ks}
K_s=(s+1)Z_sZ_{s+2}-sZ_{s+1}^2-Z_sZ_{s+1}.
\end{equation}
Let \(X(m)=1/m\) under \(\rho_s\).  By \eqref{eq:zeta-law-moment},
\begin{equation}
\label{eq:nantomah-moments}
\E_s[X]=\frac{Z_{s+1}}{Z_s},
\qquad
\E_s[X^2]=\frac{Z_{s+2}}{Z_s}.
\end{equation}
Thus
\begin{equation}
\label{eq:nantomah-moment-form}
\frac{K_s}{Z_s^2}
=(s+1)\E_s[X^2]-s\E_s[X]^2-\E_s[X]
=(s+1)\Var_s(X)+\E_s[X]^2-\E_s[X].
\end{equation}

Write \(a=\zeta(s)-1\), \(b=\zeta(s+1)-1\), and \(c=\zeta(s+2)-1\).  Expanding gives
\begin{align}
\label{eq:nantomah-lq-split}
K_s&=L_s+Q_s,&
L_s&=sa+(s+1)c-(2s+1)b,&
Q_s&=(s+1)ac-sb^2-ab.
\end{align}
The linear part is termwise positive:
\begin{equation}
\label{eq:nantomah-linear-positive}
L_s=
\sum_{m=2}^{\infty}
m^{-s}\frac{(m-1)(s(m-1)-1)}{m^2}
\ge
2^{-s}\frac{s-1}{4}.
\end{equation}
Since \(b\le a/2\) and \(c\ge0\),
\begin{equation}
\label{eq:nantomah-quadratic-correction}
Q_s\ge -sb^2-ab\ge-\frac{s+2}{4}a^2.
\end{equation}
For \(s\ge4\),
\begin{equation}
\label{eq:nantomah-tail-bound}
a=\sum_{m=2}^{\infty}m^{-s}
\le2^{-s}+\int_2^\infty x^{-s}\,dx
=2^{-s}\left(1+\frac2{s-1}\right).
\end{equation}
Equations \eqref{eq:nantomah-linear-positive}, \eqref{eq:nantomah-quadratic-correction}, and \eqref{eq:nantomah-tail-bound} imply
\begin{equation}
\label{eq:nantomah-lower-bound}
K_s\ge
\frac14\left[
2^{-s}(s-1)
-(s+2)2^{-2s}\left(1+\frac2{s-1}\right)^2
\right].
\end{equation}
It is therefore enough to show
\begin{equation}
\label{eq:nantomah-sufficient-inequality}
2^s(s-1)>(s+2)\left(1+\frac2{s-1}\right)^2.
\end{equation}
Equivalently,
\begin{equation}
\label{eq:nantomah-f-function}
F(s)=\frac{2^s(s-1)^3}{(s+2)(s+1)^2}>1.
\end{equation}
Now \(F(4)=72/25>1\), and
\begin{equation}
\label{eq:nantomah-f-monotone}
\frac{d}{ds}\log F(s)
=\log2+\frac3{s-1}-\frac1{s+2}-\frac2{s+1}>0
\qquad(s\ge4).
\end{equation}
Thus \(K_s>0\) for \(s\ge4\).

For \(s=2\),
\begin{equation}
\label{eq:nantomah-k2}
K_2=3\zeta(2)\zeta(4)-2\zeta(3)^2-\zeta(2)\zeta(3).
\end{equation}
Using \(\zeta(3)<5/4\), \(\zeta(2)=\pi^2/6\), and \(\zeta(4)=\pi^4/90\),
\begin{equation}
\label{eq:nantomah-k2-bound}
K_2>
\frac{\pi^6}{180}-\frac{5\pi^2}{24}-\frac{25}{8}
>
\frac{10415360504209}{180000000000000}>0.
\end{equation}
For \(s=3\), the estimates \(\zeta(3)>251/216\), \(\zeta(5)>8051/7776\), and \(\zeta(4)<13/12\) yield
\begin{equation}
\label{eq:nantomah-k3-bound}
K_3>
\frac{251}{216}
\left(4\cdot\frac{8051}{7776}-\frac{13}{12}\right)
-3\left(\frac{13}{12}\right)^2
=\frac{13783}{419904}>0.
\end{equation}
Hence \(K_s>0\) for all integers \(s\ge2\), proving the theorem.  This answers Nantomah's question \cite{nantomah}.
\end{proof}

\begin{theorem}[Alzer-Kwong convexity and concavity pattern for reciprocal zeta]
\label{thm:alzer-kwong}
Let \(F(x)=1/\zeta(x)\).  For every integer \(n\ge1\),
\begin{equation}
\label{eq:alzer-kwong-convex}
F''(x)>0\quad\text{on }(-4n,-4n+2),
\end{equation}
and
\begin{equation}
\label{eq:alzer-kwong-concave}
F''(x)<0\quad\text{on }(-4n-2,-4n).
\end{equation}
\end{theorem}

\begin{proof}
Put \(x=-u\), where \(u>2\).  The functional equation gives
\begin{equation}
\label{eq:zeta-functional-equation-negative-axis}
\zeta(-u)
=-2^{-u}\pi^{-u-1}\sin\left(\frac{\pi u}{2}\right)\Gamma(u+1)\zeta(u+1),
\end{equation}
and hence
\begin{equation}
\label{eq:reciprocal-zeta-negative-axis}
\frac1{\zeta(-u)}
=-\frac{2^u\pi^{u+1}}
{\Gamma(u+1)\zeta(u+1)\sin(\pi u/2)}.
\end{equation}
Define the positive function
\begin{equation}
\label{eq:alzer-g-function}
G(u)=
\frac{2^u\pi^{u+1}}
{\Gamma(u+1)\zeta(u+1)|\sin(\pi u/2)|}.
\end{equation}
On \(u\in(4n-2,4n)\), \(\sin(\pi u/2)<0\), so \(F(-u)=G(u)\).  On \(u\in(4n,4n+2)\), \(\sin(\pi u/2)>0\), so \(F(-u)=-G(u)\).  Since second derivatives are unchanged by \(x=-u\), it remains to prove \(G''(u)>0\).

Let \(\ell(u)=\log G(u)\).  Then
\begin{equation}
\label{eq:alzer-log-g}
\ell(u)=u\log2+(u+1)\log\pi-\log\Gamma(u+1)-\log\zeta(u+1)
-\log\left|\sin\left(\frac{\pi u}{2}\right)\right|.
\end{equation}
Differentiating twice gives
\begin{equation}
\label{eq:alzer-log-g-second}
\ell''(u)
=-\psi'(u+1)-(\log\zeta)''(u+1)
+\frac{\pi^2}{4}\csc^2\left(\frac{\pi u}{2}\right).
\end{equation}
With \(s=u+1>3\),
\begin{equation}
\label{eq:trigamma-bound}
\psi'(s)=\sum_{k=0}^{\infty}\frac1{(s+k)^2}
<
\int_{s-1}^{\infty}\frac{dt}{t^2}
=\frac1{s-1}\le\frac12.
\end{equation}
Combining \eqref{eq:uniform-curvature-bound}, \eqref{eq:trigamma-bound}, and \(\csc^2(\pi u/2)\ge1\), we get
\begin{equation}
\label{eq:alzer-log-convexity}
\ell''(u)>
\frac{\pi^2}{4}-\frac12-\frac13
=\frac{\pi^2}{4}-\frac56>0.
\end{equation}
Since \(G>0\),
\begin{equation}
\label{eq:alzer-g-second-positive}
G''(u)=G(u)(\ell''(u)+\ell'(u)^2)>0.
\end{equation}
The two signs in \eqref{eq:alzer-kwong-convex} and \eqref{eq:alzer-kwong-concave} now follow from \(F(-u)=G(u)\) on \(u\in(4n-2,4n)\) and \(F(-u)=-G(u)\) on \(u\in(4n,4n+2)\).  This proves the Alzer--Kwong sign pattern \cite{alzer-kwong}.
\end{proof}

\begin{theorem}[Generalized Holder inequality for Gamma zeta]
\label{thm:sroysang-holder}
Let \(m\ge2\), \(r\ge1\), \(x_1,\ldots,x_m>1\), and \(p_1,\ldots,p_m>1\) satisfy
\begin{equation}
\label{eq:sroysang-holder-condition}
\sum_{i=1}^{m}\frac1{p_i}=\frac1r.
\end{equation}
Define
\begin{equation}
\label{eq:sroysang-t}
T=r\sum_{i=1}^{m}\frac{x_i}{p_i}.
\end{equation}
Then \(T>1\) and
\begin{equation}
\label{eq:sroysang-generalized-holder}
\zeta(T)\le
\frac{\prod_{i=1}^{m}(\Gamma(x_i)\zeta(x_i))^{r/p_i}}{\Gamma(T)}.
\end{equation}
\end{theorem}

\begin{proof}
Set
\begin{equation}
\label{eq:sroysang-lambda}
\lambda_i=\frac{r}{p_i}.
\end{equation}
Then \(\lambda_i>0\), \(\sum_i\lambda_i=1\), and \(T=\sum_i\lambda_i x_i>1\).  Define
\begin{equation}
\label{eq:sroysang-q}
q_i=\frac{p_i}{r}.
\end{equation}
Then \(\sum_i1/q_i=1\), and the assumptions imply \(q_i>1\) for all \(i\); otherwise one term \(1/p_i\) would already be at least \(1/r\), leaving no room for the remaining positive terms in \eqref{eq:sroysang-holder-condition}.

For \(t>0\), let
\begin{equation}
\label{eq:sroysang-fi}
f_i(t)=\left(\frac{t^{x_i-1}}{e^t-1}\right)^{1/p_i}.
\end{equation}
Then
\begin{equation}
\label{eq:sroysang-product}
\left(\prod_{i=1}^{m}f_i(t)\right)^r
=\frac{t^{r\sum_i(x_i-1)/p_i}}{e^t-1}
=\frac{t^{T-1}}{e^t-1}.
\end{equation}
Consequently,
\begin{equation}
\label{eq:sroysang-lr-norm}
\left\|\prod_{i=1}^{m}f_i\right\|_{L^r(0,\infty)}^r
=\int_0^\infty\frac{t^{T-1}}{e^t-1}\,dt
=\Gamma(T)\zeta(T).
\end{equation}
Generalized Holder gives
\begin{equation}
\label{eq:sroysang-holder-step}
\left\|\prod_{i=1}^{m}f_i\right\|_{L^r}
\le
\prod_{i=1}^{m}\|f_i\|_{L^{p_i}}.
\end{equation}
Moreover,
\begin{equation}
\label{eq:sroysang-pi-norm}
\|f_i\|_{L^{p_i}}^{p_i}
=\int_0^\infty\frac{t^{x_i-1}}{e^t-1}\,dt
=\Gamma(x_i)\zeta(x_i).
\end{equation}
Combining \eqref{eq:sroysang-lr-norm}, \eqref{eq:sroysang-holder-step}, and \eqref{eq:sroysang-pi-norm}, then raising to the power \(r\), gives
\begin{equation}
\label{eq:sroysang-predivide}
\Gamma(T)\zeta(T)
\le
\prod_{i=1}^{m}(\Gamma(x_i)\zeta(x_i))^{r/p_i}.
\end{equation}
Division by \(\Gamma(T)\) proves \eqref{eq:sroysang-generalized-holder}.  This solves the generalized Holder problem in Sroysang's notation, where his \(\xi(s)\) is \(\zeta(s)\) for \(s>1\) \cite{sroysang}.
\end{proof}

\begin{remark}[Free-energy interpretation of the Sroysang inequality]
The same result follows from log-convexity of \(M(s)=\Gamma(s)\zeta(s)\).  With \(\lambda_i=r/p_i\) and \(T=\sum_i\lambda_ix_i\), Jensen's inequality for \(\log M\) gives
\begin{equation}
\label{eq:sroysang-jensen}
M(T)\le\prod_{i=1}^{m}M(x_i)^{\lambda_i}
=\prod_{i=1}^{m}(\Gamma(x_i)\zeta(x_i))^{r/p_i},
\end{equation}
which is equivalent to \eqref{eq:sroysang-generalized-holder}.
\end{remark}

\section{Conclusion}

The zeta-law normalization gives a single language for several identities and inequalities involving \(\zeta\).  At the structural level, the paper proves the successor entropy formula \eqref{eq:successor-entropy-series} and its modular resolution \eqref{eq:successor-entropy-modular-resolution}, the Euler-score decomposition \eqref{eq:euler-score-decomposition}, and the Mellin-Planck variance identity \eqref{eq:mellin-planck-log-convexity}.  These results are not only formal translations: they supply direct estimates for external problems.  The Nantomah positivity question follows from a moment split and tail bounds; the Alzer--Kwong reciprocal-zeta curvature pattern follows from the functional equation and a positive-axis curvature estimate; and the Sroysang generalized Holder inequality follows from the Planck-kernel integral representation.  The remaining useful direction is to sharpen the modular entropy side, especially the prime-modulus functionals in \eqref{eq:prime-functional-limit}, where additive successor entropy is encoded by finite Dirichlet \(L\)-value data.

\begin{thebibliography}{9}

\bibitem{nantomah}
K. Nantomah,
\emph{Open Problem on Riemann Zeta Function},
ResearchGate problem note, October 2024.
\url{https://www.researchgate.net/publication/384676538_Open_Problem_on_Riemann_Zeta_Function}.

\bibitem{sroysang}
B. Sroysang,
\emph{Two Inequalities for the Riemann Zeta Functions},
Mathematica Aeterna 3, no. 1 (2013), 21--24.
\url{https://arastirmax.com/en/system/files/dergiler/135290/makaleler/3/1/arastirmax-two-inequalities-riemann-zeta-functions.pdf}.

\bibitem{alzer-kwong}
H. Alzer and M. K. Kwong,
\emph{On the concavity and convexity of \(1/\zeta\)},
International Journal of Number Theory 21, no. 8 (2025), 1825--1835.
\url{https://doi.org/10.1142/S1793042125500897}.

\bibitem{apostol}
T. M. Apostol,
\emph{Introduction to Analytic Number Theory},
Undergraduate Texts in Mathematics, Springer, 1976.
\url{https://doi.org/10.1007/978-1-4757-5579-4}.

\bibitem{davenport}
H. Davenport,
\emph{Multiplicative Number Theory},
3rd ed., revised by H. L. Montgomery, Graduate Texts in Mathematics, vol. 74, Springer, 2000.
\url{https://doi.org/10.1007/978-1-4757-5927-3}.

\bibitem{titchmarsh}
E. C. Titchmarsh,
\emph{The Theory of the Riemann Zeta-function},
2nd ed., revised by D. R. Heath-Brown, Oxford University Press, 1986.
\url{https://global.oup.com/academic/product/the-theory-of-the-riemann-zeta-function-9780198533696}.

\bibitem{cover-thomas}
T. M. Cover and J. A. Thomas,
\emph{Elements of Information Theory},
2nd ed., Wiley-Interscience, 2006.
\url{https://doi.org/10.1002/047174882X}.

\end{thebibliography}

\end{document}
